\subsubsection{Equações de Lagrange e estrutura das forças de inércia}
\label{sec:dinAcoplada_equacoes_Lagrange}

As equações de Lagrange para o sistema espaçonave--\gls{mr}, em termos das coordenadas generalizadas $\vec q\,$ e de suas derivadas, podem ser escritas componente a componente conforme \cite{Wilde2018}:

\begin{equation}
\frac{d}{dt}\!\left(\frac{\partial T}{\partial \dot q_k}\right)
-\frac{\partial T}{\partial q_k}
=
\tau_k,
\qquad
k = 1,\dots,n_g,
\label{eq:lagrange_componentes}
\end{equation}

onde $T$ é a energia cinética total, $\tau_k$ é a força generalizada associada à coordenada $q_k$ e $n_g$ é o número de graus de liberdade considerados na dinâmica acoplada (espaçonave e \gls{mr}).
%%%%%%%%%%%%%%%%%%%%%%%%%%%%%%%%%%%%%%%%%%%%%%%%%%%%%%%%%%%%%%%%%%%%%%%%%%%%%%%%%%%%%%%%%%%%%%%

Partindo da forma quadrática da energia cinética já estabelecida em
\eqref{eq:T_matriz_corrigida}, em que $M(\vec q\,)$ é a matriz de inércia
generalizada do conjunto espaçonave--\gls{mr}, pode-se reescrever $T$ no
espaço escalar das coordenadas generalizadas.

\begin{equation}
T(\vec q,\dot{\vec q\, })
=
\frac{1}{2}
\sum_{k=1}^{n_g}\sum_{j=1}^{n_g}
m_{kj}(\vec q\, )\,\dot q_k \dot q_j.
\label{eq:T_expandida_mkj}
\end{equation}

Sendo assim, a derivada parcial de $T$ em relação à velocidade generalizada
$\dot q_k$ é obtida de forma explícita, e ao reescrever $T$ na forma

\begin{equation}
T(\vec q,\dot{\vec q\, })
=
\frac{1}{2}
\sum_{j=1}^{n_g}\sum_{l=1}^{n_g}
m_{jl}(\vec q)\,\dot q_j \dot q_l,
\end{equation}

tem-se

\begin{equation}
\begin{aligned}
\frac{\partial T}{\partial \dot q_k}
&=
\frac{\partial}{\partial \dot q_k}
\biggl[
\frac{1}{2}
\sum_{j=1}^{n_g}\sum_{l=1}^{n_g}
m_{jl}(\vec q)\,\dot q_j \dot q_l
\biggr]
\\
&=
\frac{1}{2}
\sum_{l=1}^{n_g}
m_{kl}(\vec q)\,\dot q_l
+
\frac{1}{2}
\sum_{j=1}^{n_g}
m_{jk}(\vec q)\,\dot q_j
\\
&=
\sum_{j=1}^{n_g}
m_{kj}(\vec q)\,\dot q_j
\end{aligned}
\label{eq:dT_dqdot_intermediaria}
\end{equation}

onde, no último passo, utiliza-se a simetria da matriz de inércia generalizada,
$m_{jk}(\vec q) = m_{kj}(\vec q)$ \cite{Geradin2004}, em seguida, ao aplicar a derivada temporal total a \eqref{eq:dT_dqdot_intermediaria}, obtém-se

\begin{equation}
\begin{aligned}
\frac{d}{dt}\!\left(\frac{\partial T}{\partial \dot q_k}\right)
&=
\frac{d}{dt}
\biggl(
\sum_{j=1}^{n_g}
m_{kj}(\vec q)\,\dot q_j
\biggr)
\\
&=
\sum_{j=1}^{n_g}
\bigl(
\dot m_{kj}(\vec q)\,\dot q_j
+
m_{kj}(\vec q)\,\ddot q_j
\bigr),
\end{aligned}
\end{equation}

onde

\begin{equation}
\dot m_{kj}(\vec q)
=
\sum_{l=1}^{n_g}
\frac{\partial m_{kj}}{\partial q_l}\,\dot q_l.
\end{equation}

Substituindo esta expressão em $\dot m_{kj}(\vec q)$ na relação anterior e
reorganizando os termos, obtém-se uma forma compacta para o primeiro termo
à esquerda de \eqref{eq:lagrange_componentes}, associado a cada coordenada
$q_k$ do sistema espaçonave--\gls{mr}.
Por sua vez, o segundo termo de \eqref{eq:lagrange_componentes} envolve a
derivada parcial de $T$ em relação à coordenada generalizada $q_k$. Usando
novamente a forma escalar em \eqref{eq:T_expandida_mkj}, tem-se

\begin{equation}
\begin{aligned}
\frac{\partial T}{\partial q_k}
&=
\frac{\partial}{\partial q_k}
\biggl[
\frac{1}{2}
\sum_{j=1}^{n_g}\sum_{l=1}^{n_g}
m_{jl}(\vec q)\,\dot q_j \dot q_l
\biggr]
\\
&=
\frac{1}{2}
\sum_{j=1}^{n_g}\sum_{l=1}^{n_g}
\frac{\partial m_{jl}}{\partial q_k}\,\dot q_j \dot q_l,
\end{aligned}
\end{equation}

onde as velocidades generalizadas $\dot q_j$ e $\dot q_l$ são tratadas como constantes na derivação em
relação a $q_k$. Essa expressão fornece o segundo termo cinético das equações de Lagrange em
\eqref{eq:lagrange_componentes} para o sistema espaçonave--\gls{mr}.


%%%%%%%%%%%%%%%%%%%%%%%%%%%%%%%%%%%%%%%%%%%%%%%%%%%%%%%%%%%%%%%%%%%%%%%%%%%%%%%%%%%%%%%%%%%%%%%

Combinando as expressões obtidas para
$\dfrac{d}{dt}\!\left(\dfrac{\partial T}{\partial \dot q_k}\right)$
e para $\dfrac{\partial T}{\partial q_k}$, pode-se escrever as equações
de Lagrange na forma escalar

\begin{equation}
\sum_{j=1}^{n_g}
m_{kj}(\vec q\,)\,\ddot q_j
+
\sum_{j=1}^{n_g}\sum_{l=1}^{n_g}
\bigg(
\frac{\partial m_{kj}}{\partial q_l}
-
\frac{1}{2}\frac{\partial m_{lj}}{\partial q_k}
\bigg)\,\dot q_l \dot q_j
=
\tau_k - g_k(\vec q\,),
\qquad
k = 1,\dots,n_g,
\label{eq:lagrange_componentes_expandida}
\end{equation}

onde $q_k$ e $q_l$ são as componentes da coordenada generalizada
$\vec q$, e $g_k(\vec q\,)$ é a $k$-ésima componente do vetor
$\vec G(\vec q\,)$ associado às forças generalizadas derivadas do
potencial configuracional $V(\vec q\,)$, com
$g_k(\vec q\,) = \dfrac{\partial V}{\partial q_k}$ \cite{Geradin2004}.
Ao agrupar os termos proporcionais às acelerações $\ddot q_j$ na matriz
$M(\vec q\,)$, os termos quadráticos em $\dot q_j$ e $\dot q_l$ na
matriz $C(\vec q\,,\dot{\vec q\,})$ e as contribuições do potencial no
vetor $\vec G(\vec q\,)$, a expressão
\eqref{eq:lagrange_componentes_expandida} é organizada em forma
matricial \cite{Geradin2004}:

\begin{equation}
M(\vec q\,)\,\ddot{\vec q\,}
+
C(\vec q\,,\dot{\vec q\,})\,\dot{\vec q\,}
=
\vec\tau
-
\vec G(\vec q\,),
\label{eq:dinamica_geral_com_G}
\end{equation}

onde o segundo termo agrupa as contribuições centrífugas e de Coriolis,
e a matriz $C(\vec q\,,\dot{\vec q\,})$ possui elementos dados por
\eqref{eq:coeficientes_C_geradin}.

Conforme discutido na Seção~\ref{sec:dinAcoplada_energia_cinetica},
no modelo adotado para o conjunto espaçonave--\gls{mr} em ambiente de
microgravidade assume-se $V(\vec q\,) \equiv 0$, de modo que
$\vec G(\vec q\,) = \vec 0$ e a Lagrangiana reduz-se a
$L(\vec q\,,\dot{\vec q\,}) = T(\vec q\,,\dot{\vec q\,})$
\cite{wie2008space,murray1994mathematical}.
Nesse caso, as equações de movimento tomam a forma padrão

\begin{equation}
M(\vec q\,)\,\ddot{\vec q\,}
+
C(\vec q\,,\dot{\vec q\,})\,\dot{\vec q\,}
=
\vec\tau,
\label{eq:dinamica_geral}
\end{equation}

onde $M(\vec q\,)$ é a matriz de inércia generalizada construída a partir
dos jacobianos espaciais e dos operadores de inércia da base e dos elos,
conforme \eqref{eq:M_total_formula_corrigida}, e
$C(\vec q\,,\dot{\vec q\,})\,\dot{\vec q\,}$ representa o vetor de
forças centrífugas--Coriolis \cite{featherstone2008,lynch2017modern}.

A organização das coordenadas generalizadas segue a definição
\eqref{eq:q_generalizado_base_mr}, isto é,

\begin{equation}
\vec q
=
\begin{bmatrix}
\vec\eta_B \\
\vec\theta
\end{bmatrix},
\qquad
\dot{\vec q}
=
\begin{bmatrix}
\dot{\vec\eta}_B \\
\dot{\vec\theta}
\end{bmatrix},
\end{equation}

onde $\vec\eta_B$ reúne os parâmetros de atitude da base e $\vec\theta$
contém as coordenadas articulares do manipulador.
De forma compatível, o vetor de forças generalizadas é divido como

\begin{equation}
\vec\tau
=
\begin{bmatrix}
\vec\tau_{\text{att}} \\
\vec\tau_{\text{juntas}}
\end{bmatrix},
\label{eq:tau_estruturado_final}
\end{equation}

onde $\vec\tau_{\text{att}}$ é o vetor de torques aplicados à atitude da
espaçonave, gerados pelo controlador de atitude, e
$\vec\tau_{\text{juntas}}$ é o vetor de torques aplicados nos atuadores
das juntas do \gls{mr}, gerados pelos controladores articulares.
Essa estrutura deixa explícito que, no modelo considerado, $\vec\tau$
agrupa apenas as ações de controle internas ao conjunto
espaçonave--\gls{mr}, enquanto as contribuições derivadas de potenciais
externos (como torques gravitacionais) são negligenciadas ao se adotar
$V(\vec q\,) = 0$.

%%%%%%%%%%%%%%%%%%%%%3333333333333%%%%%%%%%%%%%%%%%%%%%%%%%%%%%%%%%%%%%%%%

Sendo assim, é conveniente explicitar a Equação~\eqref{eq:dinamica_geral}
para o conjunto espaçonave--\gls{mr} por meio do particionamento,
conforme \cite{featherstone2008,Wilde2018}, da matriz de inércia e da
matriz de Coriolis:

\begin{equation}
M(\vec q\,) =
\begin{bmatrix}
M_{BB}(\vec q\,) & M_{BM}(\vec q\,) \\
M_{MB}(\vec q\,) & M_{MM}(\vec q\,)
\end{bmatrix},
\qquad
C(\vec q\,,\dot{\vec q\,}) =
\begin{bmatrix}
C_{BB}(\vec q\,,\dot{\vec q\,}) & C_{BM}(\vec q\,,\dot{\vec q\,}) \\
C_{MB}(\vec q\,,\dot{\vec q\,}) & C_{MM}(\vec q\,,\dot{\vec q\,})
\end{bmatrix},
\label{eq:particionamento_M_C}
\end{equation}

onde o índice $B$ está associado às contribuições da base física e o
índice $M$ às contribuições do manipulador robótico.
Ao aplicar esse particionamento na Equação~\eqref{eq:dinamica_geral},
a dinâmica da atitude da base pode ser escrita na forma

\begin{equation}
M_{BB}(\vec q\,)\,\ddot{\vec q}_B
+
M_{BM}(\vec q\,)\,\ddot{\vec q}_M
+
C_{BB}(\vec q\,,\dot{\vec q\,})\,\dot{\vec q}_B
+
C_{BM}(\vec q\,,\dot{\vec q\,})\,\dot{\vec q}_M
=
\vec\tau_B,
\label{eq:dinAcoplada_base_compacta}
\end{equation}

onde $\vec q_B$ é o vetor de coordenadas generalizadas associado à
atitude da base, $\vec q_M$ reúne as coordenadas das juntas do
manipulador e $\vec\tau_B$ é o vetor de forças generalizadas aplicadas
à base.
Para explicitar o papel das reações inerciais geradas pelo
manipulador, decompõe-se o esforço generalizado da base em

\begin{equation}
\vec\tau_B
=
\vec\tau_{\text{att}}
+
\vec\tau_{\text{reacao}},
\end{equation}

onde $\vec\tau_{\text{att}}$ representa o torque de controle aplicado
pelos atuadores de atitude da espaçonave e $\vec\tau_{\text{reacao}}$
representa o torque interno transmitido ao corpo da base pela
movimentação do manipulador, isto é, a reação de acoplamento dinâmica
entre base e \gls{mr}. 